\documentclass[12pt, twoside]{article}
%\documentclass[12pt, twoside]{article}
\usepackage[letterpaper, margin=1in, headsep=0.2in]{geometry}
\setlength{\headheight}{0.6in}
%\usepackage[english]{babel}
\usepackage[utf8]{inputenc}
\usepackage{microtype}
\usepackage{amsmath}
\usepackage{amssymb}
%\usepackage{amsfonts}
\usepackage[nomessages]{fp} %\FPeval{\var-name}{2*sin(pi/6)}
\usepackage{siunitx} %units in math. eg 20\milli\meter
\usepackage{yhmath} % for arcs, overparenth command
\usepackage{tikz} %graphics
\usetikzlibrary{quotes, angles, arrows, arrows.meta}
\usepackage{pgfplots}
\usepgfplotslibrary{fillbetween}
\pgfplotsset{compat=1.16}
\usepackage{graphicx} %consider setting \graphicspath{{images/}}
\usepackage{parskip} %no paragraph indent
\usepackage{enumitem}
\usepackage{multicol}
\usepackage{venndiagram}

\usepackage{fancyhdr}
\pagestyle{fancy}
\fancyhf{}
\renewcommand{\headrulewidth}{0pt} % disable the underline of the header
\raggedbottom
\hfuzz=2mm %suppresses overfull box warnings

\usepackage{hyperref}

\fancyhead[LO]{La Scuola d'Italia / Dr.\ Huson / IB Math 12 \\* 18 November 2026}
\fancyhead[RO]{\fbox{\begin{minipage}{6cm}\footnotesize Nickname:\\[0.5cm]\end{minipage}}}
\fancyfoot[C]{\thepage}

\begin{document}
\subsubsection*{3.7 Classwork: Applications of Periodic Functions}

All angles are in radians: set your GDC to radian mode. Give answers
correct to three significant figures.

\begin{enumerate}[itemsep=0.15cm]

%--- Problem 1: Ferris wheel h(t) = d - a cos(bt): a and d from the wheel, b from the period, h(2), first time at 40 m, time above 40 m per revolution [9 marks] ---
%--- Source: Original (AA SL 2.11, 3.7); easier than 3-6Q version (a) Problem 4: no phase shift, model form given ---
\item[1.]
\begin{minipage}[t]{0.6\linewidth}
A Ferris wheel has a diameter of $40$ metres. Its centre is $25$ metres
above the ground, and it turns once every $12$ minutes. A seat starts at
the lowest point when $t = 0$.

\smallskip
The height of the seat above the ground, in metres, $t$ minutes later is
modelled by
\[
h(t) = d - a\cos(bt), \qquad a,\ b,\ d > 0.
\]
\end{minipage}%
\begin{minipage}[t]{0.4\linewidth}
\vspace{0pt}
\begin{flushright}
\begin{tikzpicture}[scale=0.85]
  \draw[thick] (-2.4,0) -- (2.4,0);
  \draw[thick] (0,2.5) circle (2);
  \fill (0,2.5) circle (1.5pt);
  \fill (0,0.5) circle (2.5pt);
  \node[left, font=\small] at (-0.12,0.5) {seat};
  \draw[<->] (-2.2,0) -- (-2.2,2.5) node[midway, left, font=\small] {$25$};
  \draw[dashed] (-2.3,2.5) -- (0,2.5);
  \draw[<->] (2.15,0.5) -- (2.15,4.5) node[midway, right, font=\small] {$40$};
  \draw[dashed] (0,4.5) -- (2.3,4.5) (0,0.5) -- (2.3,0.5);
\end{tikzpicture}
\end{flushright}
\end{minipage}

\begin{enumerate}[itemsep=0.1cm]
  \item Find the value of $a$ and the value of $d$.
  \item Find the value of $b$.
  \item Find the height of the seat after $2$ minutes.
  \item Find the first time at which the seat is $40$ metres above the
  ground.
  \item Find the length of time, during one turn of the wheel, for which the
  seat is more than $40$ metres above the ground.
\end{enumerate}

\begin{tikzpicture}
  \draw (-0.6,0) rectangle (14.6,10.1);
  \foreach \y in {1,2,...,9} \draw[dotted] (0,\y)--(14.1,\y);
\end{tikzpicture}

\newpage

%--- Problem 2: Tides D(t) = a sin(b(t - c)) + d from high tide 8.4 m at 04:00 and low tide 1.6 m at 10:15: a, d; period and b; c; depth at 13:00; time the harbour is too shallow between high tides [11 marks] ---
%--- Source: Original (AA SL 2.11, 3.7); parallel in difficulty to 3-6Q version (a) Problem 4 (sine model, c from the maximum, time below a level) ---
\item[2.]

The depth of water in a harbour, in metres, $t$ hours after midnight is
modelled by
\[
D(t) = a\sin\bigl(b(t - c)\bigr) + d, \qquad a,\ b,\ d > 0, \quad
0 < c < 4.
\]
On one day, high tide is at 04:00, when the depth is $8.4$ metres. The next
low tide is at 10:15, when the depth is $1.6$ metres.

\begin{enumerate}[itemsep=0.1cm]
  \item Find the value of $a$ and the value of $d$.
  \item Find the period of the model, in hours, and hence find the value
  of $b$.
  \item Find the value of $c$.
  \item Find the depth of the water at 13:00.
  \item A fishing boat can enter the harbour only when the depth of water is
  at least $3.5$ metres. Find the length of time between the high tides at
  04:00 and 16:30 during which the boat cannot enter the harbour.
\end{enumerate}

\begin{tikzpicture}
  \draw (-0.6,0) rectangle (14.6,13.1);
  \foreach \y in {1,2,...,12} \draw[dotted] (0,\y)--(14.1,\y);
\end{tikzpicture}

\end{enumerate}
\end{document}
